\section{国家的基本制度}
\begin{enumerate}
  \item 基础性、方向性、全局性、根本性的重大问题。
  \item 政治。政体、国体。
        \begin{enumerate}
          \item 政治协商（政治性组织，93 入宪）、爱国统战。
          \item 坚持\zd{党的领导}，坚持四项基本原则。
          \item \zd{人}民代表\zd{大}会制度。权利来自人民，人大行使权利，人民监督人大，人大监督其他。
        \end{enumerate}
  \item 经济。市场重要，鼓励、支持、引导、监督、管理。
        \begin{enumerate}
          \item \zd{生产资料}公有制。社会主义\zd{市场经济}。
          \item 农村和城市郊区的土地原则上集体所有。
          \item \zd{绝对集体所有}。宅基地、自留地、自留山、水塘、水库。
          \item 社会主义的公共财产、公民\zd{合法}的私有财产神圣不可侵犯。
        \end{enumerate}
  \item 文化。\zd{科}学事业、\zd{教}育制度、\zd{文}学艺术和其他文化事业、公民\zd{道}德教育。送分点。
  \item 社会。送分的热点。
        \begin{enumerate}
          \item 五位一体（政治、经济、文化、社会、生态），新发展格局。
          \item \zd{社}会\zd{保}障、\zd{劳动}保障、\zd{医}疗\zd{卫}生、计划\zd{生育}、\zd{人才}培养、社会\zd{秩序及安全}维护（国防、国安）。
        \end{enumerate}
  \item 选举法。性价比很低。
        \begin{enumerate}
          \item 基本原则。\zd{中国国籍，年满 18，既没发疯，也没被抓；如果被抓，判前看检法，判后看附加}。
          \item \zd{直接选举、间接选举并用}。最重要考点。
                \begin{enumerate}
                  \item 县乡直接选举。选委会主持，人常任命。
                  \item 省、自治区、直辖市、设区的市、自治州 \yjt 下一级人大选举产生。本级人常主持。
                  \item 特区的选举全人常主持。选举会议成员 10+ 提名，联名提名不能超过应选人数。差额选举。
                \end{enumerate}
          \item 代表名额分配。
          \item 选举程序。
                \begin{enumerate}
                  \item 直接选举。
                        \begin{enumerate}
                          \item 选区划分、选民登记（长期有效）。
                          \item 选民名单 20 公布，5 日提意见，3 天处理，投票前 5 日起诉，投票前判决。
                          \item 推荐候选人。单独、联合（10+）。选举上一级人大代表，不限于本级人大代表。
                          \item 公布候选人名单。15 非正式，7 正式。
                          \item 双过半当选。过半投票程序有效，过半选票当选。
                          \item 委托投票。书面 + 代理人\zd{要有选民资格}（最多接 3 人）+ 选委会同意 。
                          \item 罢免。全过半。
                        \end{enumerate}
                  \item 间接选举。全过半。
                \end{enumerate}
          \item 代表资格审查。热点。政审。条件 + 程序 + 其他违法行为。
          \item 罢免。谁产生谁罢免。向人常提罢免案，过半数 + 公告；主席团，1/10 代表，1/5 常委提罢免案，报上一级人常。
          \item 辞职。直选找本单位，过半数 + 公告；间选来源地，过半数 + 通知。
          \item 补选。人大/人常补选\zd{个别}。可等额可差额。
          \item 费用列入财政预算，\zd{国库}（政府在人民银行开设的公共账户）开支。
        \end{enumerate}
  \item 国家结构形式。央地关系。单一制。
        \begin{enumerate}
          \item \zd{一宪一体制，一籍央统地}。
          \item 民族和谐，港澳回归。
          \item 区划类型。普通区划、民族\zd{自治}地方区划、特别行政区。《行政区划管理条例》
          \item \zd{全国人大建省设特（设立、撤销、更名），乡镇全找省级政府，改动县以上要找国务院}。
        \end{enumerate}
  \item 国家标志。新锐考点。
        \begin{enumerate}
          \item 国歌。\zd{人大政协会议开闭，宣誓升旗庆表纪}。
          \item 国旗。
                \begin{enumerate}
                  \item 每日升挂国旗的机构 \hz 场所。党（党委、纪委）、监察、国务院驻港澳的机构；民主党派人民团体。
                  \item 可以悬挂 \yjt 改为了应当悬挂。
                \end{enumerate}
          \item 国徽。\sjt 标准版本。中国人大网和国务院网。
                \begin{enumerate}
                  \item 五星照耀天安门，周围谷穗和齿轮。
                  \item 各级人常；政府；监察委；国务院驻港澳。线上线下同步挂。
                \end{enumerate}
        \end{enumerate}
  \item 民族区域自治。送分。非常简洁，明了。
        \begin{enumerate}
          \item 单一制下的自治。
          \item 满足双重条件才能享受。民族 + 所在地。
          \item 民族自治机关。仅指自治区、州、县的\zd{人大和政府}。
          \item \sjt\, + 一个民族乡。人大和政府领导必须是少数民族公民。
          \item 自治权。自治立法、变通停止执行权。
          \item 税收优惠但报省级政府批准。
          \item 国务院批准可以搞外贸、边贸；公安部队。
          \item 自治区、自治州可以搞对外交流。没有自治县。
        \end{enumerate}
  \item 特别行政区。1 分。比较啰嗦。
        \begin{enumerate}
          \item 高度自治。包括司法。
          \item 对外。不以国家为单位，非政治领域。
          \item 一国两制。舞照跳，马照跑，基本生活方式保持 50 年不变。
          \item 永久性居民。
          \item 涉及央管事务、央地关系的法律解释 \yjt 全人常。
          \item 公职人员就职宣誓。法定程序，必经条件。
          \item 特首。40 岁，住满 20 年，中国公民。选举委员会。188 提名，5 界齐全，均摊 15 人。
          \item 行政机关。司局处署厅，澳门没有署，特首来提名，国务院任免。
          \item 行政职权。\zd{土地自然资源}管理、货币\zd{金融}政策、\zd{民}用\zd{航}空运输、\zd{出入境}管理。
          \item 立法机关。立法会。
                \begin{enumerate}
                  \item 立法权。
                  \item 特首签署，公布生效 \yjt 报人常备案 \yjt \zd{只能发回} \yjt 立即失效，没有溯及力。
                  \item 立法会议员。都要政审，不允许申诉。\zd{香港立法会，八成自己人}。
                  \item 财政权。预算行政长官签署，报国务院备案。
                  \item 监督权。\zd{听取}施政报告、\zd{质询}政府工作、\zd{辩论}公共利益、\zd{弹劾}行政长官。报国务院决定。
                \end{enumerate}
          \item 司法机关。选项考点。
                \begin{enumerate}
                  \item 香港没有检察院。
                  \item 独委会推荐，行政长官任命。报全人常。
                \end{enumerate}
          \item 法律制度。一条红线，两种打法。
                \begin{enumerate}
                  \item 基本法。小宪法。
                  \item 保留法。
                  \item 附件三。香港国安法。涉外难以管辖，港府无效执法、现实猥亵重大。
                  \item 律政司才能提起公诉。
                \end{enumerate}
        \end{enumerate}
  \item 基层群众自治。\zd{村委会}、居委会（组织法）。
        \begin{enumerate}
          \item \zd{组织独立，事务自决}。不是领导、被领导关系；不需要报批（只需备案）。
          \item 村民会议。年满 18，发疯被抓。
                \begin{enumerate}
                  \item 满 18 过半数 \hz\, 2/3 户头可以开会。
                  \item 村民代表会议。
                  \item 村委会成员。自生自灭。不得制定、委派、撤换。
                  \item 选任。双过半。
                  \item 村民选举委员会。乡级人民政府监督。
                \end{enumerate}
          \item 村务公开。
          \item 全村公决。
          \item 经济责任审计。
          \item 归约备案。
          \item 法院撤销 \hz 政府\zd{责令改正} \hz 村民会议撤销 \hz 依法罢免村委会成员 \hz 连续两次民主评议不称职。
        \end{enumerate}
\end{enumerate}
